\section{Methodology}
\label{sec:methodology}

\subsection{From piecewise windows to a global GP fit}

The original HPS GPR studies explored piecewise analysis windows and explicit sideband
optimization at every mass point. Those studies were useful for establishing the basic
viability of GP-based background interpolation, but they also exposed shortcomings that
became increasingly important once smooth upper-limit curves and combined results were
required. Window-by-window fitting generated limit structures that were visibly less
smooth than the present global fit, increased sensitivity to local fit boundaries, and
complicated the interpretation of shared-$\eps^2$ combinations. For these reasons, the
current implementation performs a single global fit over the full histogram range of
each dataset and excludes only a narrow mass-dependent training window around the test
hypothesis.

\subsection{Observable, scan grid, and blinding}

The input observable is the reconstructed $e^+e^-$ invariant-mass histogram for each
dataset. The default scan step in the package is
\begin{equation}
\Delta m = \SI{1}{MeV},
\end{equation}
and the blind window for a mass hypothesis $m_0$ is centered on the detector
resolution,
\begin{equation}
\mathcal{B}(m_0) =
\left[m_0 - n_{\mathrm{blind}}\sigma_d(m_0),\,
      m_0 + n_{\mathrm{blind}}\sigma_d(m_0)\right].
\end{equation}
For the completed studies summarized in this note, the blind half-width is
$n_{\mathrm{blind}}=1.64$, corresponding to the 95\% CL workflow requested for the
current validation campaign. The same value is also used for the nominal GP training
exclusion window, although the package exposes this choice independently through
\texttt{gp\_train\_exclude\_nsigma} and \texttt{inj\_train\_exclude\_nsigma}.

\subsection{Log-space preprocessing and the choice \texorpdfstring{$\alpha=1/y$}{alpha=1/y}}

Following the recent HEP GPR background-modeling literature
\cite{BarrLiu2025,Gandrakota2023}, the package fits the GP in transformed coordinates:
\begin{equation}
x = \log m,
\qquad
z = \log y,
\end{equation}
where $y$ is the observed bin count. The package then assigns a heteroscedastic
variance term $\alpha_i$ to each training point and adds it to the kernel diagonal.
The current analysis adopts
\begin{equation}
\alpha_i =
\begin{cases}
1/y_i, & y_i > 0, \\
1, & y_i = 0,
\end{cases}
\label{eq:alpha-oney}
\end{equation}
with an explicit zero-count floor \texttt{pre\_zero\_alpha = 1.0}.

This is one of the key deliberate departures from the Barr--Liu prescription. The
motivation is that \texttt{GaussianProcessRegressor} expects a variance-like diagonal
contribution, and the package therefore treats $\alpha_i$ as the variance associated
with the transformed observation rather than as a regularization heuristic. In that
sense the present workflow follows the statistical logic of Poisson-like counting data
while removing the extra $\log y$ regularization factor that entered the earlier
CMS-motivated discussion in the rough drafts.

\subsection{Kernel choice and CMS-motivated background modeling}
\label{sec:kernel-choice}

The baseline kernel in the current package is a single
ConstantKernel$\times$RBF model,
\begin{equation}
k(x,x') = C \exp\!\left[-\frac{(x-x')^2}{2\ell^2}\right],
\end{equation}
with broad amplitude bounds and a resolution-scaled RBF length scale.
This choice is consistent with the smooth-background prior used in contemporary
CMS-oriented GPR signal-extraction studies \cite{Gandrakota2023} and with the
log-space HEP background-modeling strategy emphasized in Ref.~\cite{BarrLiu2025}. The
current HPS workflow intentionally keeps the kernel simple: the amplitude parameter $C$
controls the overall covariance scale, while $\ell$ controls the characteristic
bending scale of the background. That division is well matched to a smooth trident
spectrum in which the dominant statistical question is how rapidly the background is
allowed to vary relative to the detector resolution.

The earlier repository prototypes retained more elaborate kernel experiments, including
piecewise analyses and ensemble-style diagnostics. Those alternatives remain useful for
historical context, but they are not part of the current production baseline. The note
therefore documents the single-kernel Constant$\times$RBF model as the authoritative
choice.

\subsection{Resolution-scaled length-scale parameterization and bounds}
\label{sec:length-scale}

The GP is trained in $x=\log m$, so detector resolution must be translated into
log-mass units. The package defines
\begin{equation}
\sigma_x(m) = \log\!\left(m+\sigma_d(m)\right) - \log m
\simeq \frac{\sigma_d(m)}{m}
\qquad (\sigma_d \ll m),
\end{equation}
and uses this as the local physical scale that sets the RBF bounds. In the default
\path{resolution_scaled_local} policy, the length scale satisfies
\begin{equation}
\ell_{\min}(m) = k_{\min}\sigma_x(m),
\qquad
\ell_{\max}(m) = k_{\max}\sigma_x(m),
\end{equation}
with a baseline lower factor $k_{\min}=0.5$.

The physics-facing ceiling choice adopted in this note is
\begin{equation}
k_{\max} =
\begin{cases}
8, & \text{2015 and 2016}, \\
9, & \text{2021}.
\end{cases}
\end{equation}
This is the analyst-chosen interpretation of the length-scale-bound studies stored in
the repository's length-scale plots and rough drafts, and it reflects the decision that
the upper bound should be tethered to a physically meaningful multiple of the detector
resolution rather than left free-floating. A free-floating ceiling gave the GP too much
freedom to smooth over signal-like structure and could create the false appearance of
improved fit quality. Conversely, a ceiling set too low forces the optimizer to sit on
the bound and degrades the background description. The adopted 8$\sigma$ and 9$\sigma$
ceilings are therefore presented as stability choices rather than as arbitrary tunings.

The helper study YAMLs in this repository still contain an older 2021 upper factor of
6 in some validation configs. That value is treated here as a legacy workflow setting,
not as the final analysis prescription.

In addition to the local scaling, the package applies two practical guards already
visible in the current configuration:
\begin{itemize}
\item a dataset-statistic floor on $\ell_{\max}$, implemented through
\path{kernel_ls_local_hi_floor_mode = dataset_stat}, which prevents
unphysically small ceilings at masses where the local $\sigma_x(m)$ is unusually
small;
\item an optional cap at a fixed fraction of the training xrange
(\path{kernel_ls_local_hi_cap_xrange_frac = 0.25} in the current study
configs), which prevents the optimizer from turning the GP into an effectively global
polynomial interpolant.
\end{itemize}

\begin{figure}[H]
\centering
\begin{subfigure}[t]{0.48\linewidth}
  \centering
  \graphicorplaceholder{\linewidth}{kernel_lengthscale_figs/ls_hi_bounds_logm_overlay.png}
  \caption{Upper bound in log-mass units.}
\end{subfigure}
\hfill
\begin{subfigure}[t]{0.48\linewidth}
  \centering
  \graphicorplaceholder{\linewidth}{kernel_lengthscale_figs/dm_equiv_from_lhi_overlay.png}
  \caption{Equivalent mass scale implied by $\ell_{\max}$.}
\end{subfigure}
\caption{Representative length-scale-bound studies to be copied into the Overleaf
\texttt{kernel\_lengthscale\_figs/} directory. These figures motivate the use of
resolution-tied ceilings rather than unconstrained smoothing scales.}
\label{fig:lengthscale-bounds}
\end{figure}

\subsection{Blind-window prediction in count space}

For each mass point, the package trains the GP only on the bins outside the training
exclusion window and predicts the latent mean and covariance in the blinded bins. Since
the fit is performed in log-count space, the output is mapped back to count space
through the standard lognormal moment relations:
\begin{align}
\mu_i &= \exp\!\left(\mu^{(\log)}_i + \frac{1}{2}\Sigma^{(\log)}_{ii}\right), \\
\mathrm{Cov}(y_i,y_j)
&= \mu_i\mu_j\left[\exp\!\left(\Sigma^{(\log)}_{ij}\right)-1\right].
\end{align}
The resulting blind-window mean vector $\bm{b}$ and covariance matrix $\bm{C}$ are the
central statistical objects used by both the extraction fit and the \CLs\ limit.

\subsection{Signal model and conversion to \texorpdfstring{$\eps^2$}{epsilon2}}

The signal template is a bin-integrated Gaussian centered at the test mass $m_0$ with
width $\sigma_d(m_0)$:
\begin{equation}
w_{d,i}(m_0) = \int_{e_i}^{e_{i+1}}
\frac{1}{\sqrt{2\pi}\sigma_d(m_0)}
\exp\!\left[-\frac{(m-m_0)^2}{2\sigma_d(m_0)^2}\right]\,dm,
\end{equation}
renormalized so that $\sum_i w_{d,i}=1$ over the blinded bins. The signal-yield
parameter $A$ is the total number of signal events in the blinded region, so the
per-bin expectation is $A\,w_i$.

The package then converts between $A$ and $\eps^2$ via
\begin{equation}
\eps^2 =
\frac{2\alpha_{\mathrm{EM}}\,A}
{3\pi\,m\,f_{\mathrm{rad},d}(m)\,\rho_d(m)},
\label{eq:eps2-conversion}
\end{equation}
where $\rho_d(m)$ is the local integral density computed directly from the observed
histogram in the interval $m\pm2\sigma_d(m)$. Equation~\eqref{eq:eps2-conversion}
matches the current implementation in \texttt{hps\_gpr.conversion}.

\subsection{Local significance and single-dataset signal extraction}

The present local-significance baseline uses a profiled Gaussian multi-bin
likelihood-ratio test. In the blinded window, the Poisson mean in bin $i$ is
\begin{equation}
\lambda_i(A,\bm{\theta}) = b_i + (L\bm{\theta})_i + A\,w_i,
\qquad
\bm{C}=LL^\top,
\end{equation}
with Gaussian nuisance parameters $\bm{\theta}\sim\mathcal{N}(0,I)$ encoding correlated
background deformations in the blind window. The signal-extraction fit profiles the
likelihood over $(A,\bm{\theta})$ and returns $(\hat A,\sigma_A)$, allowing $\hat A<0$
for closure studies and enforcing the physical boundary $A\ge0$ for discovery and
limit-setting test statistics.

The local discovery significance is computed from the profiled likelihood-ratio test
statistic $q_0=-2\ln\lambda(A=0)$, with the package reporting
\begin{equation}
Z_{\mathrm{local}} = \sqrt{q_0},
\qquad
p_0 = 1-\Phi(Z_{\mathrm{local}}).
\end{equation}
This replaces the older sum-count approximation that appears in the rough drafts.

\subsection{\texorpdfstring{\CLs}{CLs} upper limits}

For a fixed signal hypothesis, the package constructs the log-likelihood ratio
\begin{equation}
\ln Q(A) = -\sum_i s_i(A) + \sum_i n_i \ln\!\left(1+\frac{s_i(A)}{b_i}\right),
\end{equation}
and computes \CLs\ either with an asymptotic approximation or with toy Monte Carlo.
In asymptotic mode the variance of $\ln Q$ includes both Poisson counting fluctuations
and the GP background covariance contribution $c^\top\bm{C}c$, where
$c_i=\ln(1+s_i/b_i)$. The upper limit $A_{95}$ is the smallest amplitude satisfying
\begin{equation}
\mathrm{CL}_s(A_{95}) = 0.05.
\end{equation}
Expected bands are derived from background-only pseudoexperiments. The public-facing
plots written by the package provide the median expected limit together with the
central $\pm1\sigma$ and $\pm2\sigma$ bands in both signal-yield and $\eps^2$ space.

\subsection{Shared-\texorpdfstring{$\eps^2$}{epsilon2} combination across datasets}
\label{sec:combination}

At a common mass hypothesis $m$, each active dataset contributes its own blinded-bin
count vector, GP background mean, covariance, and Gaussian signal template. The package
defines a dataset-specific conversion factor
\begin{equation}
K_d(m) \equiv A_d(\eps^2=1)
= \frac{3\pi\,m\,f_{\mathrm{rad},d}(m)\,\rho_d(m)}{2\alpha_{\mathrm{EM}}},
\end{equation}
so that the signal contribution for unit $\eps^2$ is $K_d(m)w_d(m)$. The combined fit
then concatenates all active datasets:
\begin{equation}
\bm{n} = (\bm{n}_1,\bm{n}_2,\ldots),
\quad
\bm{b} = (\bm{b}_1,\bm{b}_2,\ldots),
\quad
\bm{s}_{\mathrm{unit}} = (K_1\bm{w}_1,K_2\bm{w}_2,\ldots),
\end{equation}
with a block-diagonal background covariance
\begin{equation}
\bm{C}_{\mathrm{comb}} = \mathrm{blockdiag}(\bm{C}_1,\bm{C}_2,\ldots).
\end{equation}
This construction corresponds to the standard HEP product likelihood for statistically
independent channels with a shared parameter of interest \cite{Cowan2011,HistFactory}.

The combined upper limit is therefore set directly on the shared coupling,
\begin{equation}
\eps^2_{95}(m) = \min\{\eps^2 : \mathrm{CL}_s(\eps^2)=0.05\},
\end{equation}
while the combined local significance is evaluated from the profiled Gaussian
likelihood ratio built on the concatenated signal template per unit $\eps^2$.

\subsection{Global significance and the look-elsewhere placeholder}

The package currently reports local/global $p_0$ and $Z$ curves using a
Sidak-corrected effective-trials estimate derived from the scan spacing, the blind
window width, and the detector resolution \cite{GrossVitells2010}. This is an
appropriate placeholder for validation and bookkeeping, but the global-significance
statements in the present study should be treated as provisional until the rerun with
the corrected global-$p$ treatment is complete. Accordingly, the final note retains a
dedicated look-elsewhere section while deferring any definitive global-significance
claim.

\analysisplaceholder{\textbf{LEE placeholder.} Replace this subsection with the final
global-significance treatment once the dedicated rerun is complete. The present
repository-level summary plots are useful for layout and bookkeeping, but the final
physics note should quote only the corrected global $p_0$ and corresponding $Z$
curves.}

\subsection{Signal absorption and alternative background-profile treatments}

The production baseline documented here is the additive Gaussian-profiled treatment of
blind-window background uncertainty. That choice is motivated by the fact that the GP
predicts an additive mean vector and covariance in the blinded bins, and the current
toy suite is built around the same object. Earlier drafts also explored a fixed
background model and a log-normal profiling variant. Those alternatives are not the
default package behavior, but they remain relevant as future-validation tools,
particularly if additional toy studies confirm nontrivial signal absorption.

Two unresolved questions deserve explicit tracking:
\begin{itemize}
\item whether reoptimizing the GP length scale on every toy in the full procedural
pseudoexperiments partially absorbs injected signal;
\item whether a fixed-background or Poisson-rate treatment would better preserve signal
linearity in cases where the Gaussian-profile baseline is overly permissive.
\end{itemize}
For that reason, the present note documents the Gaussian-profile method as the current
baseline while reserving a dedicated robustness program for alternative nuisance
treatments rather than presenting them as equally mature production choices.
